\chapter{Preface}

\section{Who I am}

I am Jasper from \TeX{}.SX. I spent years learning Ti\textit{k}Z and 3D,
under the mentorship of helpful people on the internet, and eventually
formed some ideas of my own. I am especially grateful to early illustrations
from the internet which inspired me to learn to make my own.

Aside from 3D and ti\textit{k}Z, I enjoy animals and nature. I spent 
7 years volunteering in therapeutic horse riding, as well as many summers 
working in animal care. One of my favourite experiences was getting to
work with at a bird rehabilitation centre---with skunks, squirrels, owls, 
a bald eagle, and yes, even ducks. Most of my work was with horses though.

\section{What this manual is about}

This is the manual for lua-tikz3dtools, a package which is designed to 
teach 3D philosophy in a restricted yet still capable evironment.
lua-tikz3dtools is a set of tools for drawing projectively transformed,
and properly occluded simplicial scenes. This is a lot of jargon, and we 
will soon learn what to do to learn it and then how to use it.

Simplicial scenes are scenes composed of simplices, which is plural for 
simplex. A simplex in 3D is any element of the collection of all points,
line segments, and triangles in 3D. In lua-tikz3dtools, we typically 
focus on scenes composed of line segments and triangles.

Occlusion is the problem of determining what should be seen first if 
there is visual overlap between the simplices. Projective transformations
are the collection of all translations, rotations, shears, reflections,
and perspective transformations in 3D. It's a superset of the affine 
transformations, which are themselves a superset of the linear 
transformations.

All the things which I think are prerequisite to this book are in the 
appendix, and I recommend that readers start there for their journey.
If you dive into chapter 1 immediately, you will likely have difficulty
with things like affine algebra later on. It's worth it to have your ducks 
in a row before going off on an adventure.

\section{The scope of this manual}

This manual is all about learning my paradigms for projectively
transformed and properly occluded simplicial illustrations. This is 
just a slice of the broader 3D graphics scene, albeit a pedagogically 
valuable slice. We will not cover things like pixelated, spline-based, 
or implicit 3D illustrations. We do hope though that our affine 
algebraic---an important projective subset---principles will aid 
you in synthesizing your own ideas.




